% ------------------------------------------------------------------
% Data section subsection: data engineering, panel construction,
% and key variables. Generated from the RSpeciale pipeline
% (scripts/01_scraping -> 02_crosswalk -> 03_build); numbers reflect
% the July 2026 panel build. % ------------------------------------------------------------------
% Data section subsection: data engineering, panel construction,
% and key variables. Generated from the RSpeciale pipeline
% (scripts/01_scraping -> 02_crosswalk -> 03_build); numbers reflect
% the July 2026 panel build. % ------------------------------------------------------------------
% Data section subsection: data engineering, panel construction,
% and key variables. Generated from the RSpeciale pipeline
% (scripts/01_scraping -> 02_crosswalk -> 03_build); numbers reflect
% the July 2026 panel build. % ------------------------------------------------------------------
% Data section subsection: data engineering, panel construction,
% and key variables. Generated from the RSpeciale pipeline
% (scripts/01_scraping -> 02_crosswalk -> 03_build); numbers reflect
% the July 2026 panel build. \input{data_construction} where needed.
% ------------------------------------------------------------------

\subsection{Data engineering and panel construction}
\label{sec:data-construction}

The analysis data set is a player--month panel assembled from two
independent sources: (i) match-level performance data from FotMob, a
widely used football statistics provider whose per-match player
ratings are algorithmically generated from Opta event data, and (ii)
contract and appearance records originating from Transfermarkt,
accessed through a proprietary player-valuation database that stores
repeated snapshots of each player's contract status over time.
Because the two sources share no common identifier, and because each
covers a different slice of the player universe, most of the data
work consists of collection, record linkage, and the temporal
alignment of contract information --- each of which is described in
turn.

\subsubsection*{Collection}

FotMob data were collected via the site's public JSON endpoints at
the fixture level: for every league--season in scope, the full
fixture list was retrieved and each completed match was queried for
its player-level statistics (rating, minutes played, goals, assists,
cards, and man-of-the-match awards) together with the match result.
Scraping was checkpointed per league with deduplication on match
identifiers, so that interrupted runs resume without double-counting.
Coverage comprises 25 European leagues for the 2024/25 and 2025/26
seasons and the five major leagues for 2022/23 and 2023/24; the
empirical sections report all results with this coverage asymmetry in
mind. Match-level observations are aggregated to the player--month
level, yielding monthly appearance counts, minutes, goal
contributions, results, and two rating aggregates (unweighted and
minutes-weighted means). Ratings outside the interval $(0,10]$ ---
rare artifacts of matches with extra time --- are discarded before
aggregation.

Contract information comes from the valuation database, which records
for each tracked player a dated snapshot of the current contract
expiry date and the date of the most recent extension, alongside a
complete Transfermarkt appearance log (match date, club, and minutes
played). The appearance log is central for two reasons: it defines
playing time for the \emph{full} player universe rather than only for
matches FotMob rated, and it provides a month-by-month club history
against which contract information can be validated.

\subsubsection*{Record linkage}

FotMob and Transfermarkt identities were linked with a purpose-built
crosswalk. Candidate pairs were generated from exact and normalized
name matches and scored on name similarity together with club
agreement, where club names were compared through a learned alias
table that maps cross-language variants of the same club (e.g.\ ``FC
K\o benhavn'' and ``FC Copenhagen'') to a common form. Pairs in which
both the player name and the (alias-normalized) club agreed were
accepted automatically; the residual ambiguous cases were reviewed by
hand. The confirmed crosswalk is enforced to be unique on the FotMob
player identifier, and the merge into the panel is performed on
player identity alone rather than on player--league pairs, so that
players who move between covered leagues retain their link.

\subsubsection*{Temporal alignment of contract information}

Contract snapshots are irregularly spaced, so each player--month is
assigned the most recent snapshot at or before that month (a
last-observation-carried-forward merge). This leaves a gap at the
start of each player's history: months in which the player already
appears in the performance data but precedes their first contract
snapshot. Because a contract's expiry date is static in the absence
of a renewal or a transfer, the first observed snapshot can be
projected \emph{backwards} onto these months. To guard against
misattribution, the backward fill is applied only when (i) the gap
between the candidate month and the first snapshot does not exceed
365 days, and (ii) the appearance log confirms that the player was at
the same club in the candidate month as at the snapshot date --- a
club change would indicate a different contract, and such months are
left unfilled. Player--months whose contract status cannot be
established by either method are excluded. The backward fill
recovers 5{,}700 player--months in the estimation sample
($\approx$5\% of it), concentrated in the opening months of the
2024/25 season, where performance coverage begins before contract
tracking for many players. All filled observations carry a source
flag, and the results are robust to their exclusion.

\subsubsection*{Panel structure and key variables}

The resulting panel contains 217{,}619 player--month observations for
7{,}656 players. The estimation sample used for rating outcomes ---
the \emph{strict} panel, which requires a valid FotMob rating and
positive rated minutes in the month --- comprises 110{,}243
player--months for 7{,}247 players; extensive-margin outcomes are
estimated on the full panel, since months without any appearance are
informative there. The key variables are defined as follows.

\begin{description}
  \item[Contract horizon.] $\mathit{DaysToExpiry}_{it}$ is the
  difference between the contract expiry date in force and the
  observation month. The treatment indicator
  $\mathit{Bosman}_{it}$ equals one when 183 days or fewer remain,
  i.e.\ when the player may legally negotiate with foreign clubs
  under the Bosman ruling; event-study specifications instead use a
  full set of months-to-expiry indicators.

  \item[Contract spells.] A new spell begins whenever a player's
  recorded expiry date jumps by more than 90 days between consecutive
  months, which captures renewals and transfers. Spell identifiers
  support fixed effects at the player--contract level, so that
  identification comes from variation \emph{within} a contract as it
  runs down, not from comparisons across differently selected
  contracts.

  \item[Extensive margin.] $\mathit{Played}_{it}$ indicates positive
  minutes in the Transfermarkt appearance log, and
  $\mathit{Minutes}_{it}$ and $\mathit{Matches}_{it}$ measure monthly
  volume. These are observed for every player--month regardless of
  FotMob coverage.

  \item[Intensive margin.] The monthly mean and minutes-weighted mean
  of FotMob match ratings, observed only conditional on playing in a
  rated match. Because the rating distribution differs across
  leagues, seasons, and positions, ratings enter the regressions
  standardized within league$\times$season (and, in position-aware
  specifications, within league$\times$season$\times$position) cells;
  percentile-rank and substitution-adjusted transformations are used
  as robustness checks.

  \item[Controls and auxiliary variables.] Player age, position
  group, club identity, monthly goal contributions, and match-result
  aggregates (win share, points per match) from the rated fixtures;
  club-level financial variables (wage bills and wages-to-revenue,
  assigned from the financial year ending at the start of each
  season) support the heterogeneity analyses.
\end{description}

All processing steps are scripted and deterministic: raw scrapes are
transformed into monthly masters, linked through the crosswalk,
merged with contract histories, and written to versioned panel files
from which every table and figure in this thesis is generated.
 where needed.
% ------------------------------------------------------------------

\subsection{Data engineering and panel construction}
\label{sec:data-construction}

The analysis data set is a player--month panel assembled from two
independent sources: (i) match-level performance data from FotMob, a
widely used football statistics provider whose per-match player
ratings are algorithmically generated from Opta event data, and (ii)
contract and appearance records originating from Transfermarkt,
accessed through a proprietary player-valuation database that stores
repeated snapshots of each player's contract status over time.
Because the two sources share no common identifier, and because each
covers a different slice of the player universe, most of the data
work consists of collection, record linkage, and the temporal
alignment of contract information --- each of which is described in
turn.

\subsubsection*{Collection}

FotMob data were collected via the site's public JSON endpoints at
the fixture level: for every league--season in scope, the full
fixture list was retrieved and each completed match was queried for
its player-level statistics (rating, minutes played, goals, assists,
cards, and man-of-the-match awards) together with the match result.
Scraping was checkpointed per league with deduplication on match
identifiers, so that interrupted runs resume without double-counting.
Coverage comprises 25 European leagues for the 2024/25 and 2025/26
seasons and the five major leagues for 2022/23 and 2023/24; the
empirical sections report all results with this coverage asymmetry in
mind. Match-level observations are aggregated to the player--month
level, yielding monthly appearance counts, minutes, goal
contributions, results, and two rating aggregates (unweighted and
minutes-weighted means). Ratings outside the interval $(0,10]$ ---
rare artifacts of matches with extra time --- are discarded before
aggregation.

Contract information comes from the valuation database, which records
for each tracked player a dated snapshot of the current contract
expiry date and the date of the most recent extension, alongside a
complete Transfermarkt appearance log (match date, club, and minutes
played). The appearance log is central for two reasons: it defines
playing time for the \emph{full} player universe rather than only for
matches FotMob rated, and it provides a month-by-month club history
against which contract information can be validated.

\subsubsection*{Record linkage}

FotMob and Transfermarkt identities were linked with a purpose-built
crosswalk. Candidate pairs were generated from exact and normalized
name matches and scored on name similarity together with club
agreement, where club names were compared through a learned alias
table that maps cross-language variants of the same club (e.g.\ ``FC
K\o benhavn'' and ``FC Copenhagen'') to a common form. Pairs in which
both the player name and the (alias-normalized) club agreed were
accepted automatically; the residual ambiguous cases were reviewed by
hand. The confirmed crosswalk is enforced to be unique on the FotMob
player identifier, and the merge into the panel is performed on
player identity alone rather than on player--league pairs, so that
players who move between covered leagues retain their link.

\subsubsection*{Temporal alignment of contract information}

Contract snapshots are irregularly spaced, so each player--month is
assigned the most recent snapshot at or before that month (a
last-observation-carried-forward merge). This leaves a gap at the
start of each player's history: months in which the player already
appears in the performance data but precedes their first contract
snapshot. Because a contract's expiry date is static in the absence
of a renewal or a transfer, the first observed snapshot can be
projected \emph{backwards} onto these months. To guard against
misattribution, the backward fill is applied only when (i) the gap
between the candidate month and the first snapshot does not exceed
365 days, and (ii) the appearance log confirms that the player was at
the same club in the candidate month as at the snapshot date --- a
club change would indicate a different contract, and such months are
left unfilled. Player--months whose contract status cannot be
established by either method are excluded. The backward fill
recovers 5{,}700 player--months in the estimation sample
($\approx$5\% of it), concentrated in the opening months of the
2024/25 season, where performance coverage begins before contract
tracking for many players. All filled observations carry a source
flag, and the results are robust to their exclusion.

\subsubsection*{Panel structure and key variables}

The resulting panel contains 217{,}619 player--month observations for
7{,}656 players. The estimation sample used for rating outcomes ---
the \emph{strict} panel, which requires a valid FotMob rating and
positive rated minutes in the month --- comprises 110{,}243
player--months for 7{,}247 players; extensive-margin outcomes are
estimated on the full panel, since months without any appearance are
informative there. The key variables are defined as follows.

\begin{description}
  \item[Contract horizon.] $\mathit{DaysToExpiry}_{it}$ is the
  difference between the contract expiry date in force and the
  observation month. The treatment indicator
  $\mathit{Bosman}_{it}$ equals one when 183 days or fewer remain,
  i.e.\ when the player may legally negotiate with foreign clubs
  under the Bosman ruling; event-study specifications instead use a
  full set of months-to-expiry indicators.

  \item[Contract spells.] A new spell begins whenever a player's
  recorded expiry date jumps by more than 90 days between consecutive
  months, which captures renewals and transfers. Spell identifiers
  support fixed effects at the player--contract level, so that
  identification comes from variation \emph{within} a contract as it
  runs down, not from comparisons across differently selected
  contracts.

  \item[Extensive margin.] $\mathit{Played}_{it}$ indicates positive
  minutes in the Transfermarkt appearance log, and
  $\mathit{Minutes}_{it}$ and $\mathit{Matches}_{it}$ measure monthly
  volume. These are observed for every player--month regardless of
  FotMob coverage.

  \item[Intensive margin.] The monthly mean and minutes-weighted mean
  of FotMob match ratings, observed only conditional on playing in a
  rated match. Because the rating distribution differs across
  leagues, seasons, and positions, ratings enter the regressions
  standardized within league$\times$season (and, in position-aware
  specifications, within league$\times$season$\times$position) cells;
  percentile-rank and substitution-adjusted transformations are used
  as robustness checks.

  \item[Controls and auxiliary variables.] Player age, position
  group, club identity, monthly goal contributions, and match-result
  aggregates (win share, points per match) from the rated fixtures;
  club-level financial variables (wage bills and wages-to-revenue,
  assigned from the financial year ending at the start of each
  season) support the heterogeneity analyses.
\end{description}

All processing steps are scripted and deterministic: raw scrapes are
transformed into monthly masters, linked through the crosswalk,
merged with contract histories, and written to versioned panel files
from which every table and figure in this thesis is generated.
 where needed.
% ------------------------------------------------------------------

\subsection{Data engineering and panel construction}
\label{sec:data-construction}

The analysis data set is a player--month panel assembled from two
independent sources: (i) match-level performance data from FotMob, a
widely used football statistics provider whose per-match player
ratings are algorithmically generated from Opta event data, and (ii)
contract and appearance records originating from Transfermarkt,
accessed through a proprietary player-valuation database that stores
repeated snapshots of each player's contract status over time.
Because the two sources share no common identifier, and because each
covers a different slice of the player universe, most of the data
work consists of collection, record linkage, and the temporal
alignment of contract information --- each of which is described in
turn.

\subsubsection*{Collection}

FotMob data were collected via the site's public JSON endpoints at
the fixture level: for every league--season in scope, the full
fixture list was retrieved and each completed match was queried for
its player-level statistics (rating, minutes played, goals, assists,
cards, and man-of-the-match awards) together with the match result.
Scraping was checkpointed per league with deduplication on match
identifiers, so that interrupted runs resume without double-counting.
Coverage comprises 25 European leagues for the 2024/25 and 2025/26
seasons and the five major leagues for 2022/23 and 2023/24; the
empirical sections report all results with this coverage asymmetry in
mind. Match-level observations are aggregated to the player--month
level, yielding monthly appearance counts, minutes, goal
contributions, results, and two rating aggregates (unweighted and
minutes-weighted means). Ratings outside the interval $(0,10]$ ---
rare artifacts of matches with extra time --- are discarded before
aggregation.

Contract information comes from the valuation database, which records
for each tracked player a dated snapshot of the current contract
expiry date and the date of the most recent extension, alongside a
complete Transfermarkt appearance log (match date, club, and minutes
played). The appearance log is central for two reasons: it defines
playing time for the \emph{full} player universe rather than only for
matches FotMob rated, and it provides a month-by-month club history
against which contract information can be validated.

\subsubsection*{Record linkage}

FotMob and Transfermarkt identities were linked with a purpose-built
crosswalk. Candidate pairs were generated from exact and normalized
name matches and scored on name similarity together with club
agreement, where club names were compared through a learned alias
table that maps cross-language variants of the same club (e.g.\ ``FC
K\o benhavn'' and ``FC Copenhagen'') to a common form. Pairs in which
both the player name and the (alias-normalized) club agreed were
accepted automatically; the residual ambiguous cases were reviewed by
hand. The confirmed crosswalk is enforced to be unique on the FotMob
player identifier, and the merge into the panel is performed on
player identity alone rather than on player--league pairs, so that
players who move between covered leagues retain their link.

\subsubsection*{Temporal alignment of contract information}

Contract snapshots are irregularly spaced, so each player--month is
assigned the most recent snapshot at or before that month (a
last-observation-carried-forward merge). This leaves a gap at the
start of each player's history: months in which the player already
appears in the performance data but precedes their first contract
snapshot. Because a contract's expiry date is static in the absence
of a renewal or a transfer, the first observed snapshot can be
projected \emph{backwards} onto these months. To guard against
misattribution, the backward fill is applied only when (i) the gap
between the candidate month and the first snapshot does not exceed
365 days, and (ii) the appearance log confirms that the player was at
the same club in the candidate month as at the snapshot date --- a
club change would indicate a different contract, and such months are
left unfilled. Player--months whose contract status cannot be
established by either method are excluded. The backward fill
recovers 5{,}700 player--months in the estimation sample
($\approx$5\% of it), concentrated in the opening months of the
2024/25 season, where performance coverage begins before contract
tracking for many players. All filled observations carry a source
flag, and the results are robust to their exclusion.

\subsubsection*{Panel structure and key variables}

The resulting panel contains 217{,}619 player--month observations for
7{,}656 players. The estimation sample used for rating outcomes ---
the \emph{strict} panel, which requires a valid FotMob rating and
positive rated minutes in the month --- comprises 110{,}243
player--months for 7{,}247 players; extensive-margin outcomes are
estimated on the full panel, since months without any appearance are
informative there. The key variables are defined as follows.

\begin{description}
  \item[Contract horizon.] $\mathit{DaysToExpiry}_{it}$ is the
  difference between the contract expiry date in force and the
  observation month. The treatment indicator
  $\mathit{Bosman}_{it}$ equals one when 183 days or fewer remain,
  i.e.\ when the player may legally negotiate with foreign clubs
  under the Bosman ruling; event-study specifications instead use a
  full set of months-to-expiry indicators.

  \item[Contract spells.] A new spell begins whenever a player's
  recorded expiry date jumps by more than 90 days between consecutive
  months, which captures renewals and transfers. Spell identifiers
  support fixed effects at the player--contract level, so that
  identification comes from variation \emph{within} a contract as it
  runs down, not from comparisons across differently selected
  contracts.

  \item[Extensive margin.] $\mathit{Played}_{it}$ indicates positive
  minutes in the Transfermarkt appearance log, and
  $\mathit{Minutes}_{it}$ and $\mathit{Matches}_{it}$ measure monthly
  volume. These are observed for every player--month regardless of
  FotMob coverage.

  \item[Intensive margin.] The monthly mean and minutes-weighted mean
  of FotMob match ratings, observed only conditional on playing in a
  rated match. Because the rating distribution differs across
  leagues, seasons, and positions, ratings enter the regressions
  standardized within league$\times$season (and, in position-aware
  specifications, within league$\times$season$\times$position) cells;
  percentile-rank and substitution-adjusted transformations are used
  as robustness checks.

  \item[Controls and auxiliary variables.] Player age, position
  group, club identity, monthly goal contributions, and match-result
  aggregates (win share, points per match) from the rated fixtures;
  club-level financial variables (wage bills and wages-to-revenue,
  assigned from the financial year ending at the start of each
  season) support the heterogeneity analyses.
\end{description}

All processing steps are scripted and deterministic: raw scrapes are
transformed into monthly masters, linked through the crosswalk,
merged with contract histories, and written to versioned panel files
from which every table and figure in this thesis is generated.
 where needed.
% ------------------------------------------------------------------

\subsection{Data engineering and panel construction}
\label{sec:data-construction}

The analysis data set is a player--month panel assembled from two
independent sources: (i) match-level performance data from FotMob, a
widely used football statistics provider whose per-match player
ratings are algorithmically generated from Opta event data, and (ii)
contract and appearance records originating from Transfermarkt,
accessed through a proprietary player-valuation database that stores
repeated snapshots of each player's contract status over time.
Because the two sources share no common identifier, and because each
covers a different slice of the player universe, most of the data
work consists of collection, record linkage, and the temporal
alignment of contract information --- each of which is described in
turn.

\subsubsection*{Collection}

FotMob data were collected via the site's public JSON endpoints at
the fixture level: for every league--season in scope, the full
fixture list was retrieved and each completed match was queried for
its player-level statistics (rating, minutes played, goals, assists,
cards, and man-of-the-match awards) together with the match result.
Scraping was checkpointed per league with deduplication on match
identifiers, so that interrupted runs resume without double-counting.
Coverage comprises 25 European leagues for the 2024/25 and 2025/26
seasons and the five major leagues for 2022/23 and 2023/24; the
empirical sections report all results with this coverage asymmetry in
mind. Match-level observations are aggregated to the player--month
level, yielding monthly appearance counts, minutes, goal
contributions, results, and two rating aggregates (unweighted and
minutes-weighted means). Ratings outside the interval $(0,10]$ ---
rare artifacts of matches with extra time --- are discarded before
aggregation.

Contract information comes from the valuation database, which records
for each tracked player a dated snapshot of the current contract
expiry date and the date of the most recent extension, alongside a
complete Transfermarkt appearance log (match date, club, and minutes
played). The appearance log is central for two reasons: it defines
playing time for the \emph{full} player universe rather than only for
matches FotMob rated, and it provides a month-by-month club history
against which contract information can be validated.

\subsubsection*{Record linkage}

FotMob and Transfermarkt identities were linked with a purpose-built
crosswalk. Candidate pairs were generated from exact and normalized
name matches and scored on name similarity together with club
agreement, where club names were compared through a learned alias
table that maps cross-language variants of the same club (e.g.\ ``FC
K\o benhavn'' and ``FC Copenhagen'') to a common form. Pairs in which
both the player name and the (alias-normalized) club agreed were
accepted automatically; the residual ambiguous cases were reviewed by
hand. The confirmed crosswalk is enforced to be unique on the FotMob
player identifier, and the merge into the panel is performed on
player identity alone rather than on player--league pairs, so that
players who move between covered leagues retain their link.

\subsubsection*{Temporal alignment of contract information}

Contract snapshots are irregularly spaced, so each player--month is
assigned the most recent snapshot at or before that month (a
last-observation-carried-forward merge). This leaves a gap at the
start of each player's history: months in which the player already
appears in the performance data but precedes their first contract
snapshot. Because a contract's expiry date is static in the absence
of a renewal or a transfer, the first observed snapshot can be
projected \emph{backwards} onto these months. To guard against
misattribution, the backward fill is applied only when (i) the gap
between the candidate month and the first snapshot does not exceed
365 days, and (ii) the appearance log confirms that the player was at
the same club in the candidate month as at the snapshot date --- a
club change would indicate a different contract, and such months are
left unfilled. Player--months whose contract status cannot be
established by either method are excluded. The backward fill
recovers 5{,}700 player--months in the estimation sample
($\approx$5\% of it), concentrated in the opening months of the
2024/25 season, where performance coverage begins before contract
tracking for many players. All filled observations carry a source
flag, and the results are robust to their exclusion.

\subsubsection*{Panel structure and key variables}

The resulting panel contains 217{,}619 player--month observations for
7{,}656 players. The estimation sample used for rating outcomes ---
the \emph{strict} panel, which requires a valid FotMob rating and
positive rated minutes in the month --- comprises 110{,}243
player--months for 7{,}247 players; extensive-margin outcomes are
estimated on the full panel, since months without any appearance are
informative there. The key variables are defined as follows.

\begin{description}
  \item[Contract horizon.] $\mathit{DaysToExpiry}_{it}$ is the
  difference between the contract expiry date in force and the
  observation month. The treatment indicator
  $\mathit{Bosman}_{it}$ equals one when 183 days or fewer remain,
  i.e.\ when the player may legally negotiate with foreign clubs
  under the Bosman ruling; event-study specifications instead use a
  full set of months-to-expiry indicators.

  \item[Contract spells.] A new spell begins whenever a player's
  recorded expiry date jumps by more than 90 days between consecutive
  months, which captures renewals and transfers. Spell identifiers
  support fixed effects at the player--contract level, so that
  identification comes from variation \emph{within} a contract as it
  runs down, not from comparisons across differently selected
  contracts.

  \item[Extensive margin.] $\mathit{Played}_{it}$ indicates positive
  minutes in the Transfermarkt appearance log, and
  $\mathit{Minutes}_{it}$ and $\mathit{Matches}_{it}$ measure monthly
  volume. These are observed for every player--month regardless of
  FotMob coverage.

  \item[Intensive margin.] The monthly mean and minutes-weighted mean
  of FotMob match ratings, observed only conditional on playing in a
  rated match. Because the rating distribution differs across
  leagues, seasons, and positions, ratings enter the regressions
  standardized within league$\times$season (and, in position-aware
  specifications, within league$\times$season$\times$position) cells;
  percentile-rank and substitution-adjusted transformations are used
  as robustness checks.

  \item[Controls and auxiliary variables.] Player age, position
  group, club identity, monthly goal contributions, and match-result
  aggregates (win share, points per match) from the rated fixtures;
  club-level financial variables (wage bills and wages-to-revenue,
  assigned from the financial year ending at the start of each
  season) support the heterogeneity analyses.
\end{description}

All processing steps are scripted and deterministic: raw scrapes are
transformed into monthly masters, linked through the crosswalk,
merged with contract histories, and written to versioned panel files
from which every table and figure in this thesis is generated.
